
\documentclass{article}
\usepackage{colortbl}
\usepackage{makecell}
\usepackage{multirow}
\usepackage{supertabular}

\begin{document}

\newcounter{utterance}

\centering \large Interaction Transcript for game `codenames', experiment `frequency\_high', episode 4 with reflexion{-}Qwen3.5{-}27B.
\vspace{24pt}

{ \footnotesize  \setcounter{utterance}{1}
\setlength{\tabcolsep}{0pt}
\begin{supertabular}{c@{$\;$}|p{.15\linewidth}@{}p{.15\linewidth}p{.15\linewidth}p{.15\linewidth}p{.15\linewidth}p{.15\linewidth}}
    \# & $\;$A & \multicolumn{4}{c}{Game Master} & $\;\:$B\\
    \hline

    \theutterance \stepcounter{utterance}  
    & & & \multicolumn{2}{p{0.3\linewidth}}{
        \cellcolor[rgb]{0.9,0.9,0.9}{
            \makecell[{{p{\linewidth}}}]{
                \texttt{\tiny{[GM$|$GM]}}
                {'hidden': {'team': ['friend', 'report', 'middle'], 'innocent': ['bar', 'sir', 'plane', 'country'], 'opponent': [], 'assassin': ['kind']}, 'revealed': {'team': {'team': [], 'innocent': [], 'opponent': [], 'assassin': []}, 'opponent': {'team': [], 'innocent': [], 'opponent': [], 'assassin': []}}}
            }
        }
    }
    & & \\ \\

    \theutterance \stepcounter{utterance}  
    & & \multicolumn{4}{p{0.6\linewidth}}{
        \cellcolor[rgb]{0.9,0.9,0.9}{
            \makecell[{{p{\linewidth}}}]{
                \texttt{\tiny{[P1$\langle$GM]}}
                \texttt{You have attempted this same task before and did not succeed. These are the plans you wrote after those attempts.} \\
\texttt{Trial \#1: New plan: The previous attempt was aborted because the clue "news" is identical to a distractor word in the list, violating the rule that the clue cannot be one of the words in the lists. Additionally, the clue "news" is too closely related to the distractor "news" itself, causing confusion. Next time, I will strictly verify that my chosen clue word does not appear in the team, opponent, distractor, or kill word lists. I will also ensure the clue is semantically distinct from distractors to prevent accidental guesses. For the remaining words (friend, difference, report, dad, accident, middle, drink, age, sister), I will choose a safe clue like "story" for "report" and "accident" (avoiding "news"), or "time" for "age" and "middle", ensuring the word is not in any forbidden list.} \\
\\ 
\texttt{Trial \#2: New plan: The previous attempt was aborted because the response included extra text ("You: ") before the required format, violating the instruction to "never include any other text than is required in this form." Additionally, the clue "story" only covered two words, which is inefficient. Next time, I will output strictly the two lines "CLUE: <WORD>" and "TARGETS: <WORDS>" with no prefixes, labels, or conversational filler. I will also select a broader clue, such as "family" for "friend, dad, sister," to cover more targets while ensuring the clue word is not present in any of the forbidden lists (team, opponent, distractor, or kill words).} \\
\\ 
\texttt{Trial \#3: New plan: I failed because I adopted a conservative strategy of guessing only one word per turn (e.g., "drink", "accident", "middle"), which allowed the opponent to catch up and win before I could clear my list. Additionally, my clues were too specific, leading to "innocent" distractor guesses (like "news" for "story" or "order" for "time") that ended my turns prematurely. Next time, I will prioritize high{-}yield clues that connect multiple remaining words (e.g., using "PAPER" for "report" and "friend" if applicable, or "RELATION" for "friend" and "difference") to maximize progress per turn. I will also rigorously check that my chosen clue does not semantically trigger any remaining distractor words (like "bar", "sir", "plane", "country") or opponent words, ensuring the association is strong for my targets but weak for the forbidden lists.} \\
\\ 
\\ 
\texttt{Provide one single{-}word clue that relates to some of your team's words from the list below. You can choose to how many and to which words your clue relates to.} \\
\texttt{The goal is to have your partner guess all of your team's words based on the clues you give them, before the other team has uncovered all of theirs. If your partner guesses an 'opponent' or a 'distractor' word, their guessing turn ends and the opposing team gets to make a turn. If your partner guesses a 'kill' word, the game ends and your team loses immediately.} \\
\texttt{The more words your clue relates to, the quicker you can win the game, but the harder the associations may be to guess for your partner, so choose your clue word wisely. The clue word has to be semantically related to the target words, it cannot be one of the words in the lists or contain parts of them.} \\
\\ 
\texttt{Always give your single{-}word clue and your comma{-}separated list of related target words in the following format and make your answers as short as possible, never include any other text than is required in this form:} \\
\\ 
\texttt{CLUE: <WORD>} \\
\texttt{TARGETS: <TARGETS>} \\
\\ 
\texttt{Your team words are: friend, report, middle.} \\
\texttt{Your opponent's team words are: .} \\
\texttt{Distractor words are: bar, sir, plane, country.} \\
\texttt{The kill words are: kind.} \\
            }
        }
    }
    & \\ \\

    \theutterance \stepcounter{utterance}  
    & \multicolumn{4}{p{0.6\linewidth}}{
        \cellcolor[rgb]{0.9,0.9,0.9}{
            \makecell[{{p{\linewidth}}}]{
                \texttt{\tiny{[P1$\rangle$GM]}}
                \texttt{CLUE: CENTER} \\
\texttt{TARGETS: middle} \\
            }
        }
    }
    & & \\ \\

    \theutterance \stepcounter{utterance}  
    & & & \multicolumn{2}{p{0.3\linewidth}}{
        \cellcolor[rgb]{0.9,0.9,0.9}{
            \makecell[{{p{\linewidth}}}]{
                \texttt{\tiny{[GM$|$GM]}}
                \texttt{center} \\
            }
        }
    }
    & & \\ \\

    \theutterance \stepcounter{utterance}  
    & & & \multicolumn{2}{p{0.3\linewidth}}{
        \cellcolor[rgb]{0.9,0.9,0.9}{
            \makecell[{{p{\linewidth}}}]{
                \texttt{\tiny{[GM$|$GM]}}
                ['middle']
            }
        }
    }
    & & \\ \\

    \theutterance \stepcounter{utterance}  
    & & & \multicolumn{2}{p{0.3\linewidth}}{
        \cellcolor[rgb]{0.9,0.9,0.9}{
            \makecell[{{p{\linewidth}}}]{
                \texttt{\tiny{[GM$|$GM]}}
                {'word': 'middle', 'assignment': 'team'}
            }
        }
    }
    & & \\ \\

    \theutterance \stepcounter{utterance}  
    & & & \multicolumn{2}{p{0.3\linewidth}}{
        \cellcolor[rgb]{0.9,0.9,0.9}{
            \makecell[{{p{\linewidth}}}]{
                \texttt{\tiny{[GM$|$GM]}}
                \texttt{opponent has won} \\
            }
        }
    }
    & & \\ \\

    \theutterance \stepcounter{utterance}  
    & \multicolumn{4}{p{0.6\linewidth}}{
        \cellcolor[rgb]{0.95,0.95,0.95}{
            \makecell[{{p{\linewidth}}}]{
                \texttt{\tiny{[P2$\rightarrow$P1]}}
                \texttt{Post{-}episode Reflexion (not sent to GM):} \\
\texttt{New plan: I failed because I used a conservative strategy, providing a clue ("CENTER") that only targeted one word ("middle"), leaving "friend" and "report" unaddressed. This inefficiency allowed the opponent to win before I could clear my list. Next time, I will prioritize high{-}yield clues that connect multiple remaining words simultaneously. For the current set (friend, report, middle), I will look for a broader semantic link, such as "SCHOOL" (relating to friend, report, and middle school) or "GROUP" (relating to friend and middle), ensuring the clue is not a distractor, opponent word, or kill word. I will strictly verify that the chosen clue does not strongly associate with the distractors (bar, sir, plane, country) to prevent premature turn endings, and I will list all applicable targets in the TARGETS line to maximize progress per turn.} \\
            }
        }
    }
    & & \\ \\

\end{supertabular}
}

\end{document}
